\ifdefined\iscombined
	\lecturetitle{Formal Reasoning}
\else
\documentclass{article}
\usepackage{changepage}
\usepackage{centernot}
\usepackage{listings}
\usepackage{amsthm}
\usepackage{braket}
\usepackage{comment}
\usepackage{algorithm}
\usepackage{algpseudocode}
\usepackage{amsmath}
\usepackage{amsfonts}
\usepackage{sectsty}
\usepackage{hyperref}
\usepackage{fancyhdr}
\usepackage[top=1in,bottom=1in,left=1in,right=1in]{geometry}
\usepackage{float}
\usepackage{graphicx}
\usepackage{tabularx}
\usepackage{mathtools}
\usepackage{tikz}
\usetikzlibrary{positioning}

\date{
	\vspace{-.75em}
	{\large \today}
}

\title{
	\vspace*{-1.5em}
	{\Huge Lecture 18} \\
	\vspace{-1em}
	\rule{\textwidth/4}{.01em} \\
	\vspace{-.25em}
	{\Large Formal Reasoning}
	\vspace{-.75em}
}

\author{
	{\large Matthew Farah}
}

\hypersetup{
	colorlinks=true,
	linkcolor=blue,
	filecolor=magenta,
	urlcolor=blue,
	citecolor=blue,
	anchorcolor=blue
}

\fancyhead{}
\fancyhead[L]{\today}
\fancyhead[R]{Lecture 18 - Formal Reasoning}

\setlength{\parindent}{0cm}

\lstset{
	basicstyle=\ttfamily\small,
	frame=single,
	tabsize=4,
	breaklines=true,
	numbers=left,
	numberstyle=\tiny,
	numbersep=8pt,
	xleftmargin=1.5em,
	framexleftmargin=1.5em
}

\newcounter{example}
\renewcommand{\theexample}{\arabic{example}}

\newenvironment{example}[1][]{
	\refstepcounter{example}
	\begin{adjustwidth}{1.5em}{}
	\noindent\textbf{\ifx&#1&Example \theexample:\else Example \theexample \ (#1):\fi}
	\ignorespaces
}{
	\vspace{.5em}
	\end{adjustwidth}
}

\newcounter{subexample}
\renewcommand{\thesubexample}{\alph{subexample}}

\newenvironment{subexample}{
	\setcounter{step}{0}
	\refstepcounter{subexample}
	\vspace{1em}\par\noindent
	\begin{adjustwidth}{1.5em}{}
	\noindent\textbf{(\thesubexample)}
	\ignorespaces
}{
	\par
	\vspace{1em}
	\end{adjustwidth}
}

\newenvironment{solution}{
	\vspace{.5em}
	\par
	\begin{adjustwidth}{1.5em}{}
	\textbf{{Solution:}}
	\par
	\vspace{.5em}
}{
	\vspace{1em}
	\end{adjustwidth}
}

\newcounter{step}
\renewcommand{\thestep}{\Roman{step}}

\newenvironment{step}{
	\refstepcounter{step}
	\vspace{1em}\par\noindent
	\ignorespaces\noindent\textbf{(\thestep)}
	\ignorespaces
}{
	\par
	\vspace{1em}
}

\sectionfont{\fontsize{12}{12}\bfseries}
\subsectionfont{\fontsize{10}{12}\normalfont}
\subsubsectionfont{\fontsize{10}{12}\normalfont}

\begin{document}

\maketitle
\pagestyle{fancy}

\fi

\begin{itemize}
	\item Formal reasoning is a way by which we can analyze systems mathematically.
	\begin{itemize}
		\item Formal reasoning is necessary when approximating system behaviour is not sufficient.
		\begin{itemize}
			\item For example, in safety-critical systems.
		\end{itemize}
	\end{itemize}
	\item Hoare logic is formal reasoning technique for programs.
	\item Hoare logic consists of a  in a set of Hoare Triples.
	\begin{itemize}
		\item Hoare triples consist of:
		\begin{enumerate}
			\item A pre-condition ($\set{P}$).
			\begin{itemize}
				\item A pre-requisite condition of the system's variables for everything that follows.
			\end{itemize}
			\item A statement ($S$).
			\begin{itemize}
				\item The code being tested.
			\end{itemize}
			\item A post-condition ($R$).
			\begin{itemize}
				\item A state the program must be in after the statement has been executed, if the pre-condition was satisfied.
			\end{itemize}
		\end{enumerate}
	\end{itemize}
	\item The pre- and post-conditions are both expressed using predicate logic.
	\begin{itemize}
		\item The pre- and post-conditions can \underline{only} be expressed in terms of the program's inputs.
	\end{itemize}
	\item Hoare triples are written as:
	\begin{itemize}
		\item $\set{P}S\set{Q}$.
		\begin{itemize}
			\item Can be read as ``if $\set{P}$ holds and $S$ happens, then $\set{Q}$ must be true.''
		\end{itemize}
	\end{itemize}
\end{itemize}

\begin{example}
	Find the post conditions for the following set of Hoare Triples.
\end{example}

\begin{solution}

\end{solution}

\begin{itemize}
	\item Hoare logic is often used to reason about small parts of a code base.
	\begin{itemize}
		\item For example, methods of a system as opposed to the system itself.
	\end{itemize}
	\item The strongest post-condition $Q'$ is a post-condition that implies all other post-conditions.
	\begin{itemize}
		\item If $\set{P}S\set{Q'}$ holds, then for any post-condition $Q$, $\set{P}S\set{Q}$ also holds.
	\end{itemize}
	\item The weakest pre-condition $P'$ is a pre-condition implied by all other pre-conditions.
	\begin{itemize}
		\item If for all pre-conditions $P$, $\set{P}S\set{Q}$ holds, then $\set{P'}S\set{Q}$ also holds.
	\end{itemize}
	\item Can be thought of as:
	\begin{itemize}
		\item $Q'$ has the strictest requirements on the program's variables after execution, so if $\set{P}S\set{Q'}$ holds, $\set{P}S\set{Q}$ must also hold since $Q$ has less strict requirements.
		\item $P'$ makes the weakest assumptions about the program's variales, so if $\set{P'}S\set{Q}$ holds, $\set{P}S\set{Q}$ must also hold since $P$ has stricter requirements.
	\end{itemize}
\end{itemize}

\ifdefined\iscombined
	\newpage
\else
	\end{document}
\fi
